\documentclass[14pt,a4paper]{scrreprt}

\usepackage{cmap}
\usepackage[T1]{fontenc} 
\usepackage[utf8]{inputenc}
\usepackage[english,russian]{babel}

\usepackage{caption}
\usepackage{subcaption}

\usepackage{float}

\usepackage{enumitem}

\usepackage{graphicx}
\usepackage{multirow}


\usepackage{pgfplots}
\pgfplotsset{compat=newest}
\usepgfplotslibrary{units}

\usepackage{longtable}

\usepackage{caption}
\captionsetup{labelsep=endash}
\captionsetup[figure]{name={Рисунок}}
\captionsetup[subtable]{labelformat=simple}
\captionsetup[subfigure]{labelformat=simple}
\renewcommand{\thesubtable}{\text{Таблица }\arabic{chapter}\text{.}\arabic{table}\text{.}\arabic{subtable}\text{ --}}
\renewcommand{\thesubfigure}{\text{Рисунок }\arabic{chapter}\text{.}\arabic{figure}\text{.}\arabic{subfigure}\text{ --}}


\usepackage{textcomp}

\usepackage{amsmath}
\usepackage{amsfonts}
\usepackage{array}

\usepackage{geometry}
\geometry{left=30mm}
\geometry{right=15mm}
\geometry{top=20mm}
\geometry{bottom=20mm}
\geometry{foot=1.7cm}

\usepackage{titlesec}
\titleformat{\section}
{\normalsize\bfseries}
{\thesection}
{1em}{}
\titlespacing*{\chapter}{0pt}{-30pt}{8pt}
\titlespacing*{\section}{\parindent}{*4}{*4}
\titlespacing*{\subsection}{\parindent}{*4}{*4}

% Маркировка для списков
\def\labelitemi{$\circ$}
\def\labelitemii{$*$}


\usepackage{setspace}
\onehalfspacing % Полуторный интервал

\frenchspacing
\usepackage{indentfirst} % Красная строка

\usepackage{titlesec}
\usepackage{xcolor}
% Названия глав
\titleformat{\section}{\normalsize\textmd}{\thesection}{1em}{}

\definecolor{gray35}{gray}{0.35}

\newcommand{\hsp}{\hspace{20pt}} % длина линии в 20pt

\titleformat{\chapter}[hang]{\Huge}{\textcolor{gray35}{\thechapter.}\hsp}{0pt}{\Huge\textmd}

\titleformat{\section}{\Large}{\textcolor{gray35}\thesection}{20pt}{\Large\textmd}
\titleformat{\subsection}{\Large}{\thesubsection}{20pt}{\Large\textmd}
\titleformat{\subsubsection}{\normalfont\textmd}{}{0pt}{}

% Настройки введения

\addtocontents{toc}{\setcounter{tocdepth}{2}}
\addtocontents{toc}{\setcounter{secnumdepth}{1}}

\usepackage{tocloft,lipsum,pgffor}

\addtocontents{toc}{~\hfill\textnormal{Страница}\par}

\renewcommand{\cftpartfont}{\normalfont\textmd}

\addto\captionsrussian{\renewcommand{\contentsname}{Содержание}}
\renewcommand{\cfttoctitlefont}{\Huge\textmd}

\renewcommand{\cftchapfont}{\normalfont\normalsize}
\renewcommand{\cftsecfont}{\normalfont\normalsize}
\renewcommand{\cftsubsecfont}{\normalfont\normalsize}
\renewcommand{\cftsubsubsecfont}{\normalfont\normalsize}

\renewcommand{\cftchapleader}{\cftdotfill{\cftdotsep}}

\usepackage{listings}
\usepackage{xcolor}

\usepackage[pdftex]{hyperref} % Гиперссылки
\hypersetup{hidelinks}

% Листинги 
\usepackage{listings}

\definecolor{darkgray}{gray}{0.15}

\definecolor{teal}{rgb}{0.25,0.88,0.73}
\definecolor{gray}{rgb}{0.5,0.5,0.5}
\definecolor{b-red}{rgb}{0.88,0.25,0.41}
\definecolor{royal-blue}{rgb}{0.25,0.41,0.88}

\usepackage{listings}
\lstset{
	aboveskip=3mm,
	belowskip=3mm,
	frame=tb,
	frame=single,
	basicstyle=\footnotesize\ttfamily,
	numberstyle=\tiny\color{gray},
	keywordstyle=\color{royal-blue},
	commentstyle=\color{gray35},
	stringstyle=\color{b-red},
	numbers=left,
	numbersep=5pt,
	numberstyle=\tiny,
	showstringspaces=false, 
	captionpos=t,
	tabsize=4,
	language=Lisp
}

% какой то сложный кусок со стак эксчейндж для квадратных скобок
\makeatletter
\newenvironment{sqcases}{%
	\matrix@check\sqcases\env@sqcases
}{%
	\endarray\right.%
}
\def\env@sqcases{%
	\let\@ifnextchar\new@ifnextchar
	\left\lbrack
	\def\arraystretch{1.2}%
	\array{@{}l@{\quad}l@{}}%
}
\makeatother

% и для матриц
\makeatletter
\renewcommand*\env@matrix[1][\arraystretch]{%
	\edef\arraystretch{#1}%
	\hskip -\arraycolsep
	\let\@ifnextchar\new@ifnextchar
	\array{*\c@MaxMatrixCols c}}
\makeatother


\begin{document}

\begin{titlepage}
	
	\newgeometry{left=2cm, right=2cm, top=2.5cm, bottom=2.5cm}
	\fontsize{12pt}{12pt}\selectfont
	
	\noindent \begin{minipage}{0.13\textwidth}
		\includegraphics[width=\linewidth]{assets/bmstu-logo.png}
	\end{minipage}
	\noindent\begin{minipage}{0.85\textwidth}\centering
		\textbf{\textsc{Министерство науки и высшего образования Российской Федерации}}\\
		\textbf{\textsc{Федеральное государственное бюджетное образовательное 	учреждение высшего образования}}\\
		\textbf{\textsc{Московский государственный технический университет имени 	Н.Э.~Баумана}}\\
		\textbf{\textsc{(национальный исследовательский университет)}}\\
		\textbf{\textsc{(МГТУ им. Н.Э. Баумана)}}\\
	\end{minipage}
	
	\noindent\rule{18cm}{1.5pt}
	
	\vspace{8mm}
	
	\noindent\textnormal{ФАКУЛЬТЕТ}\hspace{5mm} \underline{\textnormal{~~~~~~~~~~~~~~~~~~«Информатика и системы управления»~~~~~~~~~~~~~~~~~~}} \newline\newline
	\textnormal{КАФЕДРА}\hspace{5mm} \underline{\textnormal{~~«Программное обеспечение ЭВМ и информационные технологии»~~}}
	\newline\newline
	\textnormal{НАПРАВЛЕНИЕ ПОДГОТОВКИ}\hspace{5mm} \underline{\textnormal{~~~~~~«09.03.04 Программная инженерия»~~~~~~~~}}
	
	\vspace{2.5cm}
	
	\begin{center}
		\Large\textbf{\textsc{ОТЧЕТ}}\\
		\Large\textbf{\textsc{ПО ЛАБОРАТОРНОЙ РАБОТЕ №18}}\\
	\end{center}
	
	\vspace{1cm}
	
	\noindent\textnormal{Название:} \hspace{15mm} \underline{\textnormal{~~~~~~Формирование и модификация списков на Prolog~~~~~~~}}\noindent
	
	\vspace{1.3cm}
	
	\noindent\textnormal{Дисциплина:} \hspace{10mm} \underline{\textnormal{~~~~~~Функциональное и логическое программирование~~~~~~}}\noindent
	
	\vspace{1.5cm}
	
	\noindent\textnormal{Студент} \hspace{17mm}
	\underline{\textnormal{{~~~~ИУ7-64Б~~~}}}
	\hspace{20mm}
	\underline{\textnormal{\hphantom{~~~~~~~~~~~~~~~~~~~~~~~~~~~}}} \hspace{14mm}
	\underline{\textnormal{~~~С. Д. Параскун~~~}}
	
	\vspace{2mm}
	\noindent\textnormal{\hphantom{Студент}} \hspace{23mm}\noindent
	\fontsize{8pt}{8pt}
	\textnormal{Группа}\hspace{40mm}\textnormal{Подпись, дата} \hspace{30mm}\noindent\textnormal{И. О. Фамилия}
	
	\vspace{0.5cm}
	
	\fontsize{12pt}{12pt}\selectfont
	\noindent\textnormal{Преподаватель} \hspace{52mm}
	\underline{\textnormal{\hphantom{~~~~~~~~~~~~~~~~~~~~~~~~~~~}}} \hspace{14mm}
	\noindent\underline{\textnormal{~~Н. Б. Толпинская~~}}
	
	\vspace{2mm}
	\noindent\textnormal{\hphantom{Студент}} \hspace{17mm}\noindent
	\fontsize{8pt}{8pt}
	\hphantom{Группа}\hspace{43mm}\textnormal{Подпись, дата} \hspace{30mm}\noindent\textnormal{И. О. Фамилия}
	
	\vspace{0.5cm}
	
	\fontsize{12pt}{12pt}\selectfont
	\noindent\textnormal{Преподаватель} \hspace{52mm}
	\underline{\textnormal{\hphantom{~~~~~~~~~~~~~~~~~~~~~~~~~~~}}} \hspace{14mm}
	\noindent\underline{\textnormal{~~Ю. В. Строганов~~}}
	
	\vspace{2mm}
	\noindent\textnormal{\hphantom{Студент}} \hspace{17mm}\noindent
	\fontsize{8pt}{8pt}
	\hphantom{Группа}\hspace{43mm}\textnormal{Подпись, дата} \hspace{30mm}\noindent\textnormal{И. О. Фамилия}
	
	\vspace{1cm}
	
	\fontsize{12pt}{12pt}\selectfont
	
	\begin{center}
		\vfill
		Москва, ~\the\year
		~г.
	\end{center}
	\restoregeometry

\end{titlepage}


\thispagestyle{empty}

\chapter{Практическое задание}

\section{Задание}

Создать базу знаний <<Предки>>, позволяющую наиболее эффективным способом (за меньшее количество шагов, что обеспечивается меньшим количеством предложений в БЗ -- правил), и используя разные варианты (примеры) одного вопроса, определить (указать, какой вопрос для какого варианта):
\begin{enumerate}
	\item по имени субъекта определить всех его бабушек (предки 2-го колена);
	\item по имени субъекта определить всех его дедушек (предки 2-го колена);
	\item по имени субъекта определить всех его бабушек и дедушек (предки 2-го колена);
	\item по имени субъекта определить его бабушку по материнской линии (предки 2-го колена);
	\item по имени субъекта определить его бабушку и дедушку по материнской линии (предки 2-го колена).
\end{enumerate}
Минимизировать количество правил и количество вариантов вопросов. Использовать конъюнктивные правила и простой вопрос.

Для одного из вариантов вопроса и конкретной БЗ составить таблицу, отражающую конкретный порядок работы системы, с объяснениями:
\begin{itemize}
	\item очередная проблема на каждом шаге и метод ее решения;
	\item каково новое текущее состояние резольвенты, как получено;
	\item какие дальнейшие действий? (Запускается ли алгоритм унификации? Каких термов? Почему этих?);
	\item вывод по результатам очередного шага и дальнейшие действия.
\end{itemize}
Т.к. резольвента хранится в виде стека, то состояние резольвенты требуется отображать в столбик: вершина -- сверху! Новый шаг надо начинать с нового состояния резольвенты!

\section{Код программы}

\begin{lstlisting}[language=Prolog]
domains

	name, gender = symbol.
	human = human(name, gender).

predicates

	isParent(human, human).
	isGrandparent(human, gender, human).

clauses

	isGrandparent(human(GrandName, GrandGender), LineGender, 
		human(Name, Gender)) :- 
			isParent(human(GrandName, GrandGender), human(TmpName, 
			LineGender)), 
			isParent(human(TmpName, LineGender), human(Name, Gender)).
	
	isParent(human("Irina", woman), human("Olga", woman)).
	isParent(human("Igor", man), human("Olga", woman)).
	isParent(human("Victor", man), human("Irina", woman)).
	isParent(human("Svetlana", woman), human("Irina", woman)).
	isParent(human("Fedor", man), human("Igor", man)).
	isParent(human("Uliya", woman), human("Igor", man)).

goal
	%isGrandparent(human(Name, woman), _, human("Olga", woman)).
	%isGrandparent(human(Name, man), _, human("Olga", woman)).
	%isGrandparent(Human, _, human("Olga", woman)).
	%isGrandparent(human(Name, woman), woman, human("Olga", woman)).
	isGrandparent(Human, woman, human("Olga", woman)).
\end{lstlisting}

\newpage
\section{Таблицы выполнения программы}

Запрос для задания 2:

isGrandparent(human(Name, man), \_, human(``Olga'', woman)).


\begin{table}[H]
	\begin{tabular}{|p{0.8cm\small}|p{4.7cm\small}|p{5.7cm\small}|p{4cm\small}|}	
		\hline
		№ шага & Состояния резольвенты и вывод: дальнейшие действия (почему?) & Для каких термов запускается алгоритм унификации: T1=T2 и каков результат (подстановка) & Дальнейшие действия: прямой ход или откат (почему и к чему приводит?)\\
		\hline
		1 & Резольвента:\begin{itemize}
			\item isParent(human(
		TmpName, \_), human(``Olga'', woman)) \item isParent(human(
		Name, man), human(TmpName, \_))\end{itemize} & isGrandparent(human(Name, man), \_, human(``Olga'', woman)) = isGrandparent(human(
		GrandName, GrandGender), LineGender, human(Name, Gender)). Унификация успешна\linebreak \textbf{Подстановка:} \{GrandName=Name, GrandGender=man, Name=``Olga'', Gender=woman\} & Прямой ход, приводит к попытке унификации isParent(human(
		TmpName, \_), human(``Olga'', woman)) и isGrandparent(human(
		GrandName, GrandGender), LineGender, human(Name, Gender))\\
		\hline 
		2 & Резольвента:\begin{itemize} \item isParent(human(
			TmpName, \_), human(``Olga'', woman)) \item isParent(human(
			Name, man), human(TmpName, \_))\end{itemize} & isParent(human(TmpName, \_), human(``Olga'', woman)) = isGrandparent(human(
		GrandName, GrandGender), LineGender, human(Name, Gender)). Унификация неуспешна & Прямой ход, переход к следующему предложению\\
		\hline
		3 & Резольвента:\begin{itemize} \item isParent(human(
			Name, man), human(``Irina'', \_))\end{itemize} & isParent(human(TmpName, \_), human(``Olga'', woman)) = isParent(human(``Irina'', woman), human(``Olga'', woman)). Унификация успешна\linebreak
		\textbf{Подстановка:} \{GrandName=Name, GrandGender=man, Name=``Olga'', Gender=woman, TempName=``Irina''\} & Прямой ход, попытка унификации isParent(human(
		Name, man), human(``Irina'', \_))\\
		\hline
	\end{tabular}
\end{table}

\begin{table}[H]
	\begin{tabular}{|p{0.8cm\small}|p{4.7cm\small}|p{5.7cm\small}|p{4cm\small}|}	
		\hline
		№ шага & Состояния резольвенты и вывод: дальнейшие действия (почему?) & Для каких термов запускается алгоритм унификации: T1=T2 и каков результат (подстановка) & Дальнейшие действия: прямой ход или откат (почему и к чему приводит?)\\
		\hline
		4 & Резольвента:\begin{itemize} \item isParent(human(
			Name, man), human(``Irina'', \_))\end{itemize} & isParent(human(
		Name, man), human(``Irina'', \_)) = isGrandparent(human(
		GrandName, GrandGender), LineGender, human(Name, Gender)). Унификация неуспешна & Прямой ход, переход к следующему предложению\\
		\hline
		5-6 & ... & Унификация неуспешна & ...\\
		\hline
		7 & Резольвента: пуста\linebreak \textbf{Вывод:} Name=``Victor'' & isParent(human(
		Name, man), human(``Irina'', \_)) = isParent(human(``Victor'', man), human(``Irina'', woman)). Унификация успешна\linebreak \textbf{Подстановка:} \{GrandName=``Victor'', GrandGender=man, Name=``Olga'', Gender=woman, TempName=``Irina''\} & Откат, переход к следующему предложению\\
		\hline
		8-9 & Резольвента: \begin{itemize} \item isParent(human(
			Name, man), human(``Irina'', \_))\end{itemize} & Унификация неуспешна & Прямой ход, переход к следующему предложению\\
		\hline
		10 & Резольвента:\begin{itemize} \item isParent(human(
			TmpName, \_), human(``Olga'', woman)) \item isParent(human(
			Name, man), human(TmpName, \_))\end{itemize} & isParent(human(
		Name, man), human(``Irina'', \_)) = isParent(human(``Uliya'', woman), human(``Igor'', man)). Унификация неуспешна & Откат, достижение конца БЗ, переход к следющему предложению относительно 3\\
		\hline
		11 & Резольвента:\begin{itemize} \item isParent(human(
			Name, man), human(``Igor'', \_))\end{itemize} & isParent(human(TmpName, \_), human(``Olga'', woman)) = isParent(human(``Igor'', man), human(``Olga'', woman)). Унификация успешна\linebreak
		\textbf{Подстановка:} \{GrandName=Name, GrandGender=man, Name=``Olga'', Gender=woman, TempName=``Igor''\} & Прямой ход, попытка унификации isParent(human(
		Name, man), human(``Igor'', \_))\\
		\hline
	\end{tabular}
\end{table}

\begin{table}[H]
	\begin{tabular}{|p{0.8cm\small}|p{4.7cm\small}|p{5.7cm\small}|p{4cm\small}|}	
		\hline
		№ шага & Состояния резольвенты и вывод: дальнейшие действия (почему?) & Для каких термов запускается алгоритм унификации: T1=T2 и каков результат (подстановка) & Дальнейшие действия: прямой ход или откат (почему и к чему приводит?)\\
		\hline
		12 & Резольвента:\begin{itemize} \item isParent(human(
			Name, man), human(``Igor'', \_))\end{itemize} & isParent(human(
		Name, man), human(``Igor'', \_)) = isGrandparent(human(
		GrandName, GrandGender), LineGender, human(Name, Gender)). Унификация неуспешна & Прямой ход, переход к следующему предложению\\
		\hline
		13-17 & ... & Унификация неуспешна & ... \\
		\hline
		18 & Резольвента: пуста\linebreak \textbf{Вывод:} Name=``Fedor'' & isParent(human(
		Name, man), human(``Igor'', \_)) = isParent(human(``Fedor'', man), human(``Igor'', man)). Унификация успешна\linebreak \textbf{Подстановка:} \{GrandName=``Fedor'', GrandGender=man, Name=``Olga'', Gender=woman, TempName=``Igor''\} & Откат, переход к следующему предложению\\
		\hline
		19 & Резольвента:\begin{itemize} \item isParent(human(
			Name, man), human(``Igor'', \_))\end{itemize} & isParent(human(
		Name, man), human(``Igor'', \_)) = isParent(human(``Uliya'', woman), human(``Igor'', man)). Унификация неуспешна & Откат, достижение конца БЗ, переход к следующему предложению относительно 11\\
		\hline
		20 & Резольвента:\begin{itemize} \item isParent(human(
			TmpName, \_), human(``Olga'', woman)) \item isParent(human(
			Name, man), human(TmpName, \_))\end{itemize} & isParent(human(TmpName, \_), human(``Olga'', woman)) = isParent(human(``Victor'', man), human(``Irina'', woman)). Унификация неуспешна & Прямой ход, переход к следующему предложению\\
		\hline
		21-23 & ... & Унификация неуспешна & Откат, достижение конца БЗ, переход к следующему предложению относительно 1\\
		\hline
		24-29 & Резольвента: пуста & isGrandparent(human(Name, man), \_, human(``Olga'', woman)) = isParent(...). Унификация неуспешна & Откат, достижение конца БЗ, завершение работы\\
		\hline
	\end{tabular}
\end{table}
		
\end{document}